\documentclass[12pt, a4paper, twoside]{article}
\usepackage{cmap}
\usepackage[utf8]{vietnam}
\usepackage{amsmath}
\usepackage{amssymb}
\usepackage{geometry}
\usepackage{listings}
\usepackage{listingsutf8}
\usepackage{xcolor}
\usepackage{fancyhdr}
\usepackage{hyperref}

\geometry{left=2cm, right=2cm, top=2cm, bottom=2cm}
\setlength{\headheight}{17.97171pt}

\lstset{
    language=Python,
    inputencoding=utf8,
    basicstyle=\ttfamily\small,
    keywordstyle=\color{blue},
    commentstyle=\color{gray},
    stringstyle=\color{red},
    breaklines=true,
    showstringspaces=false,
    numberstyle=\tiny\color{gray},
    numbers=left,
}

\pagestyle{fancy}
\fancyhf{}
\rhead{Báo cáo Toán Ứng Dụng và Thống Kê}
\lhead{Kì 4}
\cfoot{\thepage}

\title{\textbf{Báo Cáo Đồ Án:\\ \Huge Phần 1}}
\author{Tên: Lê Phạm Đăng Khiêm\\MSSV: 24120341\\Môn: Toán ứng dụng và thống kê}
\date{\\}

\begin{document}

\maketitle

\tableofcontents
\newpage
\section{Giới Thiệu Đồ Án}
Hỗ trợ Github copilot: Hỗ trợ viết báo cáo bằng latex nhanh chóng, kiểm chứng bằng numpy/sympy để giải thích kết quả trả về.
Đồ án triển khai các thuật toán đại số tuyến tính cơ bản bằng Python, gồm:
\begin{itemize}
    \item Khử Gauss (Gaussian Elimination)
    \item Tính định thức (Determinant Calculation)
    \item Tính ma trận nghịch đảo (Matrix Inverse)
    \item Tính hạng ma trận và cơ sở không gian (Rank and Basis Calculation)
    \item Các hàm hỗ trợ và định dạng hiển thị
\end{itemize}

Toàn bộ phép tính trong phần cài đặt sử dụng số thực \texttt{float}. Để ổn định số học, chương trình dùng ngưỡng \texttt{EPS = 1e-12} khi so sánh với 0.

\section{Tổng Quan Cấu Trúc Đồ Án}

\subsection{Danh Sách File}
\begin{itemize}
    \item \texttt{gaussian.py} - Hàm khử Gauss, thay thế ngược và các hàm hỗ trợ
    \item \texttt{determinant.py} - Tính định thức bằng khử Gauss
    \item \texttt{inverse.py} - Tính ma trận nghịch đảo bằng Gauss-Jordan
    \item \texttt{rank\_basis.py} - Tính hạng và cơ sở không gian
    \item \texttt{part1\_demo.ipynb} - Notebook minh họa với các bài kiểm thử
    \item \texttt{README.md} - Hướng dẫn sử dụng
    \item \texttt{report/} - Thư mục báo cáo
    \begin{itemize}
        \item \texttt{report/report.tex} - File mã nguồn LaTeX của báo cáo
        \item \texttt{report/report.pdf} - File PDF báo cáo
    \end{itemize}
\end{itemize}

\section{Thuật Toán Khử Gauss (gaussian.py)}

\subsection{Mô Tả Chung}

Tệp \texttt{gaussian.py} chứa các hàm cốt lõi cho thuật toán khử Gauss:

\subsubsection{Các Hàm Hỗ Trợ}

\paragraph{to\_float(value)}
Chuyển đổi giá trị đầu vào về kiểu \texttt{float}:
\begin{lstlisting}
def to_float(value):
    return float(value)
\end{lstlisting}

\paragraph{dinh\_dang\_hien\_thi(value)}
Định dạng số thực để hiển thị gọn (làm sạch các giá trị rất nhỏ về 0 và hiển thị số nguyên khi phù hợp).

\paragraph{in\_ma\_tran\_dep(ma\_tran)}
In ma trận dưới dạng cấu trúc bảng gọn gàng, căn lề phải

\subsubsection{Khử Ma Trận về Bậc Thang}

Hàm \texttt{khu\_ma\_tran\_ve\_bac\_thang(ma\_tran\_chuan\_hoa, b\_chuan\_hoa, eps)} thực hiện:

\begin{enumerate}
    \item Duyệt qua từng cột từ trái sang phải
    \item Tìm \textbf{pivot} - phần tử có giá trị tuyệt đối lớn nhất trong cột (Partial Pivoting)
    \item Hoán đổi hàng nếu cần để đưa pivot lên vị trí hiện tại
    \item Khử các phần tử dưới pivot bằng phép biến đổi hàng:
    %
    \[
    \text{Hàng } r := \text{Hàng } r - \frac{a_{r,c}}{a_{k,c}} \times \text{Hàng } k
    \]
    %
\end{enumerate}

Trong quá trình khử, các so sánh với 0 được thực hiện theo ngưỡng \texttt{EPS} thay vì so sánh bằng tuyệt đối để giảm ảnh hưởng sai số dấu phẩy động.
Trả về:
\begin{itemize}
    \item \texttt{ma\_tran\_sau\_khu}: Ma trận dạng bậc thang
    \item \texttt{b\_sau\_khu}: Vector b sau khi áp dụng các phép biến đổi tương ứng
    \item \texttt{so\_lan\_doi\_hang}: Số lần hoán đổi hàng
    \item \texttt{pivot\_rows}: Danh sách chỉ số hàng chứa pivot
    \item \texttt{pivot\_cols}: Danh sách chỉ số cột chứa pivot
\end{itemize}

\subsubsection{Tìm Nghiệm Hệ Phương Trình}

Hàm \texttt{tim\_nghiem\_he} xác định loại nghiệm:

\begin{enumerate}
    \item \textbf{Hệ vô nghiệm}: Khi tồn tại hàng mà tất cả phần tử bên trái = 0 nhưng phần tử bên phải $\neq 0$
    \item \textbf{Hệ có vô số nghiệm}: Khi số pivot < số cột (có biến tự do)
    \item \textbf{Hệ có nghiệm duy nhất}: Khi số pivot = số cột
\end{enumerate}

Với hệ có vô số nghiệm, hàm \texttt{nghiem\_tong\_quat\_he\_vo\_so\_nghiem} tính:
\begin{itemize}
    \item Nghiệm riêng $\mathbf{x_0}$
    \item Các vector cơ sở tương ứng với các biến tự do
    \item Biểu thức tham số cho mỗi ẩn
\end{itemize}

\subsubsection{Cách Sử Dụng}

\paragraph{gaussian\_eliminate(A, b)}
Giải hệ phương trình tuyến tính $A\mathbf{x} = \mathbf{b}$:

\begin{lstlisting}
A = [[1, 2, 0, 2], [3, 5, -1, 6], [2, 4, 1, 2], [2, 0, -7, 11]]
b = [6, 17, 12, 7]
ma_tran_sau_khu, nghiem_he, so_lan_doi_hang = gaussian_eliminate(A, b)

if nghiem_he is None:
    print("Vo nghiem")
elif len(nghiem_he) > 0 and isinstance(nghiem_he[0], str):
    print("Vo so nghiem:", nghiem_he)
else:
    print("Nghiem duy nhat:", nghiem_he)
\end{lstlisting}

\subsubsection{Ví Dụ Đầu Ra}

\textbf{Trường hợp 1: Hệ có nghiệm duy nhất}
\begin{lstlisting}
A = [[2, 1, -1], [-3, -1, 2], [-2, 1, 2]]
b = [8, -11, -3]
ma_tran_sau_khu, nghiem_he, so_lan_doi_hang = gaussian_eliminate(A, b)
print("Nghiem duy nhat:", nghiem_he)
# Output: Nghiem duy nhat: [2.0, 3.0, -1.0]
\end{lstlisting}

\textbf{Trường hợp 2: Hệ vô số nghiệm}
\begin{lstlisting}
A = [[1, 2, 3], [2, 4, 6]]
b = [5, 10]
ma_tran_sau_khu, nghiem_he, so_lan_doi_hang = gaussian_eliminate(A, b)
print("Vo so nghiem:", nghiem_he)
# Output: Vo so nghiem: ['5.0 - 2.0*t_0 - 3.0*t_1', '0.0 - 0.0*t_0 + 1.0*t_1', 't_0']
\end{lstlisting}

\textbf{Trường hợp 3: Hệ vô nghiệm}
\begin{lstlisting}
A = [[1, 2, 3], [2, 4, 6], [1, 0, 1]]
b = [5, 9, 2]
ma_tran_sau_khu, nghiem_he, so_lan_doi_hang = gaussian_eliminate(A, b)
if nghiem_he is None:
    print("Vo nghiem")
# Output: Vo nghiem
\end{lstlisting}

\paragraph{print\_gaussian\_eliminate(A, b)}
In kết quả định dạng đẹp ra màn hình

\paragraph{back\_substitution(U, c)}
Giải hệ tam giác trên $U\mathbf{x} = \mathbf{c}$ bằng thay thế ngược:
\[
x_i = \frac{c_i - \sum_{j=i+1}^{n} u_{ij} x_j}{u_{ii}}
\]

\section{Tính Định Thức (determinant.py)}

\subsection{Thuật Toán}

Định thức được tính thông qua khử Gauss:
\[
\det(A) = (-1)^{s} \times \prod_{i=1}^{n} a_{ii}
\]

trong đó:
\begin{itemize}
    \item $a_{ii}$ là các phần tử trên đường chéo của ma trận dạng bậc thang
    \item $s$ là số lần hoán đổi hàng
\end{itemize}

\textbf{Chú ý}: Nếu tồn tại hàng hoàn toàn bằng 0, định thức = 0.

\subsection{Cách Sử Dụng}

\begin{lstlisting}
A = [[2, 3], [1, 5]]
det = determinant(A)  # Tra ve float
print(f"Dinh thuc: {det}")

# Hoac in theo dinh dang dep
print_determinant(A)
\end{lstlisting}

\subsection{Các Trường Hợp Kiểm Thử}

\begin{itemize}
    \item Ma trận 2×2: $\det\begin{pmatrix} 2 & 3 \\ 1 & 5 \end{pmatrix} = 7$
    
    \item Ma trận 3×3: $\det\begin{pmatrix} 1 & 2 & 3 \\ 0 & 1 & 4 \\ 5 & 6 & 0 \end{pmatrix} = -49$
    
    \item Ma trận phụ thuộc tuyến tính: $\det\begin{pmatrix} 1 & 2 & 3 \\ 2 & 4 & 6 \\ 0 & 1 & 2 \end{pmatrix} = 0$
    
    \item Ma trận đơn vị: $\det(I) = 1$
    
    \item Ma trận đường chéo: $\det = $ tích các phần tử trên đường chéo
\end{itemize}

\section{Tính Ma Trận Nghịch Đảo (inverse.py)}

\subsection{Thuật Toán Gauss-Jordan}

Ma trận nghịch đảo được tính từ ma trận mở rộng $[A \mid I]$:

\begin{enumerate}
    \item \textbf{Khử xuống}: Biến đổi $[A \mid I]$ thành $[U \mid X]$ (U là tam giác trên)
    \item \textbf{Khử lên}: Tiếp tục khử để được $[I \mid A^{-1}]$
\end{enumerate}

Điều kiện: Ma trận A phải là vuông và khả nghịch ($\det(A) \neq 0$).

\subsection{Cách Sử Dụng}

\begin{lstlisting}
A = [[2, 3], [1, 5]]
A_inv = inverse(A)  # Tra ve danh sach cac float

if A_inv is not None:
    print_inverse(A)  # In A va A^(-1)
else:
    print("Ma tran khong kha nghich")
\end{lstlisting}

\subsection{Ứng Dụng Kiểm Thử}

Kiểm chứng: $A \times A^{-1} = I$

Ví dụ:
\[
\begin{pmatrix} 2 & 3 \\ 1 & 5 \end{pmatrix}^{-1} \approx \begin{pmatrix} 0.714286 & -0.428571 \\ -0.142857 & 0.285714 \end{pmatrix}
\]

\section{Tính Hạng và Cơ Sở (rank\_basis.py)}

\subsection{Định Nghĩa}

Với ma trận $A$ bất kỳ:

\begin{itemize}
    \item \textbf{Hạng (Rank)}: Số lượng cột pivot = số lượng hàng độc lập tuyến tính
    
    \item \textbf{Cơ sở không gian cột} (Column Space): Các cột pivot của ma trận A gốc
    
    \item \textbf{Cơ sở không gian dòng} (Row Space): Các hàng khác 0 của ma trận dạng bậc thang
    
    \item \textbf{Cơ sở không gian nghiệm} (Null Space): Các vector cơ sở ứng với những cột không phải pivot
\end{itemize}

\subsection{Thuật Toán}

\begin{enumerate}
    \item Khử ma trận A về dạng bậc thang
    \item Xác định các cột pivot và hàng pivot
    \item Tr ích xuất cơ sở từ các cột pivot
    \item Tính không gian nghiệm từ các cột tự do
\end{enumerate}

\subsection{Cách Sử Dụng}

\begin{lstlisting}
A = [[1, 2, 3], [2, 4, 6], [0, 1, 2]]
rank, col_space, row_space, null_space = rank_and_basis(A)

print(f"Hang: {rank}")
print(f"Co so khong gian cot: {col_space}")
print(f"Co so khong gian dong: {row_space}")
print(f"Co so khong gian nghiem: {null_space}")

# Hoac in theo dinh dang dep
print_rank_and_basis(A)
\end{lstlisting}

\section{Kiểm Chứng và Xác Nhận (part1\_demo.ipynb)}

\subsection{Mục Đích}

Notebook \texttt{part1\_demo.ipynb} cung cấp các bài kiểm thử toàn diện để xác nhận tính chính xác của các thuật toán.

\subsection{Nội Dung Chính}

\subsubsection{Test Gaussian Elimination}

\begin{itemize}
    \item \textbf{TEST 1}: Hệ có nghiệm duy nhất (kiểm chứng bằng NumPy)
    \item \textbf{TEST 2}: Hệ vô nghiệm
    \item \textbf{TEST 3}: Hệ có vô số nghiệm
    \item \textbf{TEST 4}: Trường hợp đặc biệt 1×1
    \item \textbf{TEST 5}: Trường hợp đặc biệt 2×4 (underdetermined)
    \item \textbf{TEST 6}: Partial pivoting (pivot đầu = 0)
\end{itemize}

\subsubsection{Test Back Substitution}

\begin{itemize}
    \item Ma trận tam giác trên với nghiệm duy nhất
    \item Ma trận tam giác với hệ vô nghiệm
    \item Ma trận tam giác với hệ vô số nghiệm
    \item Trường hợp đặc biệt 1×1
    \item Xử lý lỗi với ma trận không vuông
\end{itemize}

\subsubsection{Test Determinant}

\begin{itemize}
    \item Ma trận 2×2, 3×3, 4×4
    \item Ma trận phụ thuộc tuyến tính (định thức = 0)
    \item Ma trận đơn vị, ma trận đường chéo
\end{itemize}

\subsubsection{Test Matrix Inverse}

\begin{itemize}
    \item Ma trận 2×2, 3×3, 4×4
    \item Ma trận đơn vị
    \item Ma trận không khả nghịch
\end{itemize}

\subsubsection{Test Rank and Basis}

\begin{itemize}
    \item Ma trận vuông khả nghịch (full rank)
    \item Ma trận phụ thuộc tuyến tính (rank thiếu)
    \item Ma trận chữ nhật (underdetermined)
    \item Ma trận rank = 1
    \item Ma trận zero (rank = 0)
\end{itemize}

\subsection{Kiểm Chứng bằng NumPy và SymPy}

Notebook thực hiện kiểm chứng theo đúng các nhóm TEST trong \texttt{part1\_demo.ipynb}:

\begin{lstlisting}
# TEST A: Kiem chung nghiem Gauss voi numpy.linalg.solve
A_test = [[2, 3], [1, 5]]
b_test = [7, 11]
my_x = gaussian_eliminate(A_test, b_test)[1]
np_x = np.linalg.solve(np.array(A_test, dtype=float), np.array(b_test, dtype=float))
flag1 = np.allclose(np.array([float(v) for v in my_x]), np_x)

# TEST B: Kiem chung dinh thuc voi numpy.linalg.det
A_det = [[1, 2, 3], [0, 1, 4], [5, 6, 0]]
my_det = float(determinant(A_det))
np_det = float(np.linalg.det(np.array(A_det, dtype=float)))
flag2 = np.isclose(my_det, np_det)

# TEST C: Kiem chung nghich dao qua A * A^(-1) ~= I
A_inv_test = [[2, 3], [1, 5]]
my_inv = inverse(A_inv_test)
my_inv_np = np.array([[float(v) for v in row] for row in my_inv], dtype=float)
A_inv_np = np.array(A_inv_test, dtype=float)
I_calc = A_inv_np @ my_inv_np
flag3 = np.allclose(I_calc, np.eye(2))

# TEST D: So sanh co so KG cot/dong/nghiem voi SymPy
import sympy as sp

def ma_tran_co_so_cot(co_so, so_chieu):
    if not co_so:
        return np.zeros((so_chieu, 0), dtype=float)
    return np.array(co_so, dtype=float).T

def nam_trong_tap_sinh(v, M, tol=1e-8):
    if M.shape[1] == 0:
        return np.linalg.norm(v) <= tol
    he_so, *_ = np.linalg.lstsq(M, v, rcond=None)
    return np.linalg.norm(M @ he_so - v) <= tol

def cung_khong_gian_con(co_so_1, co_so_2, so_chieu, tol=1e-8):
    M1 = ma_tran_co_so_cot(co_so_1, so_chieu)
    M2 = ma_tran_co_so_cot(co_so_2, so_chieu)

    hang_1 = np.linalg.matrix_rank(M1, tol)
    hang_2 = np.linalg.matrix_rank(M2, tol)
    hang_ghep = np.linalg.matrix_rank(np.hstack([M1, M2]), tol)
    if not (hang_1 == hang_2 == hang_ghep):
        return False

    for j in range(M1.shape[1]):
        if not nam_trong_tap_sinh(M1[:, j], M2, tol):
            return False
    for j in range(M2.shape[1]):
        if not nam_trong_tap_sinh(M2[:, j], M1, tol):
            return False
    return True

A_basis = [[1, 2, 3], [2, 4, 6], [0, 1, 2]]
my_rank, my_col_space, my_row_space, my_null_space = rank_and_basis(A_basis)

A_sp = sp.Matrix(A_basis)
sp_col_space = [list(map(float, v)) for v in A_sp.columnspace()]
sp_row_space = [list(map(float, list(v))) for v in A_sp.rowspace()]
sp_null_space = [list(map(float, v)) for v in A_sp.nullspace()]

m = len(A_basis)
n = len(A_basis[0])
my_row_as_cols = [row[:] for row in my_row_space]
sp_row_as_cols = [row[:] for row in sp_row_space]

flag_col = cung_khong_gian_con(my_col_space, sp_col_space, so_chieu=m)
flag_row = cung_khong_gian_con(my_row_as_cols, sp_row_as_cols, so_chieu=n)
flag_null = cung_khong_gian_con(my_null_space, sp_null_space, so_chieu=n)

flag4 = flag_col and flag_row and flag_null
\end{lstlisting}

\section{Hướng Dẫn Sử Dụng (README.md)}

\subsection{Giải Thích Hàm gaussian\_eliminate}

Hàm \texttt{gaussian\_eliminate(A, b)} trả về tuple gồm 3 phần tử:

\begin{enumerate}
    \item \texttt{[0]}: Ma trận dạng bậc thang (kiểu \texttt{float})
    \item \texttt{[1]}: Nghiệm hệ
    \item \texttt{[2]}: Số lần hoán đổi hàng (int)
\end{enumerate}

\subsubsection{Ý Nghĩa của Nghiệm Hệ}

\begin{itemize}
    \item \texttt{None}: Hệ vô nghiệm
    \item \texttt{list[float]}: Hệ có nghiệm duy nhất
    \item \texttt{list[str]}: Hệ có vô số nghiệm (dạng tham số)
\end{itemize}

\subsubsection{Ví Dụ Kiểm Tra Nhanh}

\begin{lstlisting}
if nghiem_he is None:
    print("Vo nghiem")
elif len(nghiem_he) > 0 and isinstance(nghiem_he[0], str):
    print("Vo so nghiem:", nghiem_he)
else:
    print("Nghiem duy nhat:", nghiem_he)
\end{lstlisting}

\subsection{Xử Lý Sai Số Số Thực}

\begin{itemize}
    \item Dữ liệu đầu vào được chuẩn hóa về \texttt{float}
    \item Các phép kiểm tra bằng 0 sử dụng ngưỡng \texttt{EPS} (ví dụ: \texttt{abs(x) <= EPS})
    \item Cách làm này đảm bảo tốc độ cao cho ma trận kích thước lớn mà vẫn duy trì độ chính xác trong thực nghiệm
\end{itemize}
\subsection{Phân Biệt Hàm}

\begin{itemize}
    \item \texttt{gaussian\_eliminate(A, b)}: Lấy dữ liệu trả về để xử lý tiếp
    \item \texttt{print\_gaussian\_eliminate(A, b)}: In kết quả ra màn hình
\end{itemize}

\section{Ưu Điểm và Đặc Điểm Của Đồ Án}

\subsection{Hiệu Năng Và Độ Ổn Định}

\begin{itemize}
    \item Sử dụng \texttt{float} để tăng tốc cho ma trận kích thước lớn
    \item Dùng ngưỡng \texttt{EPS} để hạn chế nhiễu số học khi so sánh với 0
    \item Cân bằng tốt giữa tốc độ và độ chính xác thực nghiệm
\end{itemize}

\subsection{Partial Pivoting}

\begin{itemize}
    \item Cách chọn pivot theo giá trị tuyệt đối lớn nhất
    \item Cải thiện ổn định số 
    \item Giảm hiện tượng mất mát độ chính xác trong phép tính
\end{itemize}

\subsection{Xử Lý Nhiều Trường Hợp}

\begin{itemize}
    \item Hệ vô nghiệm, vô số nghiệm, và nghiệm duy nhất
    \item Ma trận vuông, hình chữ nhật, ma trận đặc biệt
    \item Hàng/cột toàn 0, ma trận suy biến
\end{itemize}

\subsection{Hiển Thị Rõ Ràng}

\begin{itemize}
    \item Định dạng ma trận đẹp, dễ đọc
    \item Hiển thị kết quả số thực gọn và nhất quán
    \item Đầu ra có cấu trúc và rõ ràng
\end{itemize}

\section{Các Thuật Toán Không Được Sử Dụng}

Theo yêu cầu của đồ án, các hàm/cơ chế giải sẵn không được dùng trong việc triển khai:

\begin{itemize}
    \item \texttt{numpy.linalg.solve}: Không dùng để giải hệ
    \item \texttt{numpy.linalg.inv}: Không dùng để tính ma trận nghịch đảo
    \item \texttt{scipy.linalg.qr}: Không dùng
    \item \texttt{scipy.linalg.lu}: Không dùng
    \item \texttt{sympy.linsolve}: Không dùng
    \item \texttt{Echelon form/RREF} từ SymPy: Không dùng
\end{itemize}

\textbf{Chú ý}: Các hàm thư viện trên chỉ được dùng cho \textbf{mục đích kiểm chứng} kết quả, không phải cài đặt thuật toán.

\section{Kết Luận}

Đồ án này thành công triển khai các thuật toán cơ bản của đại số tuyến tính bằng Python:

\begin{itemize}
    \item \textbf{Khử Gauss}: Giải hệ phương trình tuyến tính
    \item \textbf{Định thức}: Tính định thức thông qua khử Gauss
    \item \textbf{Ma trận nghịch đảo}: Sử dụng phương pháp Gauss-Jordan
    \item \textbf{Hạng và cơ sở}: Phân tích không gian vectơ
\end{itemize}

Tất cả các thuật toán:
\begin{enumerate}
    \item Được triển khai từ đầu (không dùng thư viện có sẵn để tính)
    \item Sử dụng \texttt{float} kết hợp \texttt{EPS} để đảm bảo ổn định số và hiệu năng
    \item Xử lý đầy đủ các trường hợp đặc biệt
    \item Được kiểm chứng bằng NumPy
\end{enumerate}

\end{document}
